% GENERATED, edit the French unit or the translation table
\mxnewpage\section*{EV3-B -- Mock exam 2}\label{EV3-B}
\addcontentsline{toc}{section}{EV3-B -- Mock exam 2}
\textbf{Paper reserved for the teacher}. Programme studied in chapters 1 to 10.\par Duration: 2 h. No calculator. Marked out of 20. 1--3 April 2027.
\subsection*{Part 1 -- Fluency (6 points)}
Twelve questions, \num{0,5} point each. Only one correct answer; no justification required. An incorrect or missing answer scores 0, with no loss of points.
\begin{enumerate}
\item Calculate $4+2\times3$.\par A. $10$\quad B. $9$\quad C. $18$\quad D. $13$
\item Calculate $\frac12+\frac13$.\par A. $\frac{1}{6}$\quad B. $\frac{5}{6}$\quad C. $\frac{5}{3}$\quad D. $\frac{2}{5}$
\item Calculate $\dfrac{2^4}{2^{2}}$.\par A. $4$\quad B. $1$\quad C. $2$\quad D. $9$
\item A value of 50 increases by \pc{20}. Final value?\par A. $70$\quad B. $60$\quad C. $10$\quad D. $40$
\item A value of 80 is obtained after a \pc{20} decrease. Initial value?\par A. $96$\quad B. $60$\quad C. $64$\quad D. $100$
\item A \pc{25} increase then a \pc{20} decrease: overall coefficient?\par A. $\frac{9}{10}$\quad B. $1$\quad C. $\frac{21}{20}$\quad D. $\frac{4}{5}$
\item The solution of $2x+4=0$ is\par A. $2$\quad B. $8$\quad C. $-2$\quad D. $-8$
\item The slope between $\pt{0;4}$ and $\pt{2;10}$ is\par A. $2$\quad B. $-3$\quad C. $3$\quad D. $6$
\item The mean of $4, 6, 8$ is\par A. $7$\quad B. $18$\quad C. $6$\quad D. $4$
\item A sector of $90^\circ$ represents what fraction of the total?\par A. $\frac{1}{2}$\quad B. $\frac{1}{3}$\quad C. $\frac{3}{4}$\quad D. $\frac{1}{4}$
\item A probability is $\dfrac{2}{5}$. Its complement is\par A. $\frac{1}{5}$\quad B. $\frac{5}{2}$\quad C. $\frac{2}{5}$\quad D. $\frac{3}{5}$
\item For $f(x)=x^2+4$, $f(-2)$ is\par A. $16$\quad B. $0$\quad C. $6$\quad D. $8$
\end{enumerate}
\mxnewpage\subsection*{Part 2 -- Independent exercises (14 points)}
Every step must be justified. Exact results are expected unless stated otherwise.
\mxneed{12}
\subsection*{Exercise 1 -- Four trials and a gain (5 points)}
A game consists of four independent trials; each trial succeeds with probability $\dfrac12$. The gross gain $G$ is 8 monetary units if all four trials succeed, and 0 otherwise. A stake may be set to make the game fair.
\begin{enumerate}
\item Calculate the probability of receiving 8. \hfill(1 pt)
\item Calculate the expectation and variance of G. \hfill(2 pts)
\item Determine the fair stake; explain whether the expectation is a possible gain. \hfill(2 pts)
\end{enumerate}
\mxneed{12}
\subsection*{Exercise 2 -- A derivative of constant sign (5 points)}
Consider the function $f$ defined on $\R$ by $f(x)=(x^2+1)e^x$. We want to check whether a horizontal tangent necessarily corresponds to an extremum.
\begin{enumerate}
\item Calculate and factorise $f\prime$. \hfill(1 pt)
\item Study the variation; does the horizontal tangent at -1 give an extremum? \hfill(2 pts)
\item Write the tangent at 0 then calculate $f(1)$ exactly. \hfill(2 pts)
\end{enumerate}
\mxneed{12}
\subsection*{Exercise 3 -- Projection onto a diagonal (4 points)}
In an orthonormal frame with origin $O$, consider $P\pt{2;0}$ and the line $d$ with equation $y=x$. Let $H$ be the orthogonal projection of $P$ onto $d$. The unit vector asked for at the end must be oriented from $O$ towards $H$.
\begin{enumerate}
\item Give a nonzero normal vector to the line $d$. \hfill(1 pt)
\item Determine the coordinates of $H$, then the distance $PH$. \hfill(1 pt)
\item Give the unoriented angle $\widehat{POH}$ in radians, then the coordinates of the vector $\dfrac{\vect{OH}}{OH}$. \hfill(2 pts)
\end{enumerate}
